\documentclass[11pt]{article}
\usepackage[letterpaper,margin=0.68in]{geometry}
\usepackage{amsmath,amssymb}
\usepackage{enumitem}
\usepackage{xcolor}
\usepackage{fancyhdr}
\usepackage{microtype}
\usepackage{tabularx,array,booktabs}
\usepackage{tikz}
\usepackage[most]{tcolorbox}

\definecolor{olive}{HTML}{4D5430}
\definecolor{gold}{HTML}{9A7B22}
\definecolor{pale}{HTML}{F5F1DC}
\definecolor{softgreen}{HTML}{EDF0E2}
\definecolor{softgray}{HTML}{F5F5F2}
\definecolor{ink}{HTML}{292D20}

\pagestyle{fancy}
\fancyhf{}
\lhead{\textcolor{olive}{MATH\& 107: Math in Society}}
\rhead{\textcolor{gold}{Fall 2026}}
\cfoot{\thepage}
\setlength{\headheight}{14pt}
\setlength{\parindent}{0pt}
\setlength{\parskip}{5pt}
\setlist[enumerate]{leftmargin=*,itemsep=7pt,topsep=5pt}
\setlist[itemize]{leftmargin=1.35em,itemsep=3pt,topsep=3pt}

\newtcolorbox{keybox}[1]{colback=pale,colframe=olive,boxrule=0.8pt,arc=2mm,title=\textbf{#1},fonttitle=\color{olive},left=2mm,right=2mm,top=1.5mm,bottom=1.5mm}
\newtcolorbox{examplebox}[1]{colback=softgreen,colframe=gold,boxrule=0.8pt,arc=2mm,title=\textbf{#1},fonttitle=\color{ink},left=2mm,right=2mm,top=1.5mm,bottom=1.5mm}
\newtcolorbox{refbox}[1]{colback=softgray,colframe=gold,boxrule=0.7pt,arc=2mm,title=\textbf{#1},fonttitle=\color{ink},left=2mm,right=2mm,top=1.5mm,bottom=1.5mm}

\newcommand{\summaryhead}[2]{%
  \clearpage
  {\Large\bfseries\textcolor{olive}{Module #1: #2 — Summary}}\par
  \vspace{2pt}\textcolor{gold}{\rule{\linewidth}{1.2pt}}\par\smallskip
}
\newcommand{\prequizhead}[2]{%
  \clearpage
  {\Large\bfseries\textcolor{olive}{Module #1: #2 — Prequiz}}\par
  \textit{Open-note. Show your mathematical work. A calculation and a clearly stated final answer are enough unless a sentence is specifically requested.}\par
  Name: \hspace{2.6in} Date: \hspace{1.1in}\par
  \vspace{2pt}\textcolor{gold}{\rule{\linewidth}{1.2pt}}\par\smallskip
}
\newcommand{\quizhead}[2]{%
  \clearpage
  {\Large\bfseries\textcolor{olive}{Module #1: #2 — Matching Quiz}}\par
  \textit{These problems use the same skills as the prequiz, with new values and contexts. Show enough work to make your reasoning visible.}\par
  Name: \hspace{2.6in} Date: \hspace{1.1in}\par
  \vspace{2pt}\textcolor{gold}{\rule{\linewidth}{1.2pt}}\par\smallskip
}
\newcommand{\worksmall}{\par\vspace{0.38in}}
\newcommand{\workmed}{\par\vspace{0.62in}}
\newcommand{\worklarge}{\par\vspace{0.92in}}
\newcommand{\workxl}{\par\vspace{1.22in}}
\newcommand{\choice}[1]{\(\square\)~#1\qquad}

\begin{document}
\summaryhead{7}{Cost, Revenue, and Profit}

\textbf{What this module is about.} A simple business model separates costs that exist before any item is sold from costs that increase with each unit. Revenue depends on how many units are sold and the price per unit. Profit is what remains after cost is subtracted from revenue. The break-even quantity is the point where revenue and cost are equal, and marginal profit tells us how much one additional unit changes profit.

\begin{keybox}{The three business functions}
Let \(q\) be the number of units, \(F\) fixed cost, \(v\) variable cost per unit, and \(p\) selling price per unit.
\[
C(q)=F+vq,
\qquad
R(q)=pq,
\]
\[
\pi(q)=R(q)-C(q)=(p-v)q-F.
\]
The slope of the cost line is \(v\). The slope of the revenue line is \(p\). The slope of the profit line is \(p-v\).
\end{keybox}

\begin{examplebox}{Example 1: A tutoring business}
A tutor pays \(\$180\) each week for room rental and supplies, and each tutoring session creates \(\$6\) in additional costs. The tutor charges \(\$18\) per session.
\[
C(q)=180+6q,
\qquad
R(q)=18q.
\]
Therefore
\[
\pi(q)=18q-(180+6q)=12q-180.
\]
At \(q=25\):
\[
C(25)=330,
\quad
R(25)=450,
\quad
\pi(25)=120.
\]
So the tutor earns \(\$120\) in profit at \(25\) sessions.
\end{examplebox}

\begin{keybox}{Break-even and marginal profit}
Break-even occurs when revenue equals cost, or equivalently when profit is zero:
\[
0=(p-v)q-F.
\]
Thus
\[
q^*=\frac{F}{p-v}.
\]
The quantity \(p-v\) is \textbf{marginal profit}: the amount profit changes when one additional unit is sold in this linear model.
\end{keybox}

\begin{examplebox}{Example 2: Break-even as an intersection}
For the tutoring business above,
\[
q^*=\frac{180}{18-6}=\frac{180}{12}=15.
\]
At \(15\) sessions, cost and revenue are equal. To the left of \(15\), cost is greater than revenue and profit is negative. To the right of \(15\), revenue is greater than cost and profit is positive.

Because marginal profit is \(\$12\), each additional session beyond break-even increases profit by exactly \(\$12\). Profit therefore follows an additive recursion with step size \(12\).
\end{examplebox}

\textbf{By the end of this module, you should be able to:} build \(C(q)\), \(R(q)\), and \(\pi(q)\) from a business description; evaluate them at a stated quantity; find break-even; identify marginal cost, marginal revenue, and marginal profit; and interpret the regions before and after break-even on a graph.

\prequizhead{7}{Cost, Revenue, and Profit}
\begin{refbox}{Useful formulas}
\[
C(q)=F+vq,
\qquad R(q)=pq,
\qquad \pi(q)=(p-v)q-F,
\qquad q^*=\frac{F}{p-v}.
\]
\end{refbox}

\begin{enumerate}
\item A mobile detailing service has fixed weekly costs of \(\$240\). Each car detailed creates \(\$8\) in variable costs, and customers are charged \(\$28\) per car.
\begin{enumerate}[label=(\alph*),itemsep=4pt]
\item Identify \(F\), \(v\), and \(p\).
\item Write the cost function \(C(q)\).
\item Write the revenue function \(R(q)\).
\item Write the profit function \(\pi(q)\).
\end{enumerate}
\workxl

\item Using the detailing service from Problem 1, suppose \(18\) cars are detailed in one week. Find the total cost, total revenue, and profit. State whether the business is operating at a profit or a loss.
\worklarge

\item Find the break-even number of cars for the business in Problem 1. If the break-even quantity is a whole number, explain what happens to profit at one car below break-even and one car above break-even.
\worklarge

\item A food stand has
\[
C(q)=250+4q,
\qquad
R(q)=9q,
\]
where \(q\) is the number of meals sold.
\begin{enumerate}[label=(\alph*),itemsep=4pt]
\item Find marginal cost, marginal revenue, and marginal profit.
\item Find the break-even quantity.
\item How much additional profit is created by selling \(10\) more meals after break-even?
\end{enumerate}
\worklarge

\item On a graph of cost and revenue for a business with positive marginal profit, what is true to the \emph{right} of the break-even point: \choice{cost is larger} \choice{revenue is larger} \choice{the two are always equal}. Explain how the graph tells you the sign of profit there.
\workmed
\end{enumerate}

\quizhead{7}{Cost, Revenue, and Profit}
\begin{enumerate}
\item A print shop has fixed weekly costs of \(\$240\), variable cost \(\$3\) per poster, and selling price \(\$9\) per poster. Write \(C(q)\), \(R(q)\), and \(\pi(q)\).
\worklarge

\item For the print shop in Problem 1, find cost, revenue, and profit when \(q=50\) posters are sold.
\worklarge

\item Find the exact break-even quantity for the print shop. What is the minimum whole number of posters needed for nonnegative profit?
\workmed

\item A craft business has \(F=\$300\), \(v=\$7\), and \(p=\$15\). Find marginal profit and determine how much profit increases when \(12\) additional units are sold.
\workmed

\item A business is operating at a quantity larger than its break-even quantity, and \(p>v\). Which line is higher on the graph, cost or revenue? What does that imply about profit?
\workmed
\end{enumerate}

\end{document}
